\origpage{43}{}
\noindent contains both $G_1$ and $G_2$ if and only if it topologically contains one of $H_1, \dots, H_n$. ($G$ \textit{topologically contains} $H$ if some subgraph of $G$ is isomorphic to a subdivision of $H$.) This conjecture was proved in [11] by well-quasi-ordering methods. Our object here is to give a different proof of the intertwining conjecture, one that is ``constructive'' in the sense that it yields an algorithm which, given $G_1$ and $G_2$, computes $H_1, \dots, H_n$ as above.

Poljak and Turzik [4] showed that the intertwining conjecture is (constructively) implied by another conjecture, (11.1) below, and we shall give a constructive proof of (11.1). Let $L_1$ and $L_2$ be linkages in a graph $G$. We say that $G$ is an \textit{intertwining} of $L_1$ and $L_2$ if $L'_1 \cup L'_2 = G$ for all linkages $L'_1, L'_2$ in $G$ such that $L'_i$ has the same pattern as $L_i$ ($i = 1, 2$). Let us say a graph $G$ has \textit{rank}
\begin{equation*}
|E(G)| - |V(G)| + \kappa(G),
\end{equation*}
where $\kappa(G)$ denotes the number of components of $G$. Poljak and Turzik essentially reduced proving the intertwining conjecture to proving the following.

\textbf{11.1} \textit{For every integer $p \ge 0$ there exists $r(p)$ such that every intertwining of two $p$-linkages has rank $\le r(p)$.}

We shall derive (11.1) from (1.1), and the proof will be constructive in the sense that it yields an algorithm to compute a value for $r(p)$ given a numerical value for $p$. Our thanks to A. Gupta and R. Impagliazzo, who proved (11.1) constructively for planar intertwinings (unpublished) and brought the general problem to our attention. Thanks also to R. Thomas, who collaborated with us in discovering a constructive proof of (11.1) (not exactly the one given here).

Let $(T, W)$ be a tree-decomposition of a graph $G$. If $S$ is a subtree of $T$, we define $W(S) = \bigcup (W(t) : t \in V(S))$. If $e$ is an edge of $T$, and $T_1, T_2$ are the two components of $T \setminus e$ and $t_i$ is the end of $e$ in $T_i$ ($i = 1, 2$), then $(W(T_1), W(T_2))$ is a separation of $G$\fix{, provided that no edge of $G$ belongs to $W(t)$ for two distinct $t \in V(T)$,} and $W(T_1) \cap W(T_2) = W(t_1) \cap W(t_2)$.\fixnote{fix:11-separation}{As printed, the first claim is false: a separation $(A, B)$ has $E(A \cap B) = \emptyset$, and the definition of a tree-decomposition allows two of the graphs $W(t)$ to share an edge. For example, let $G$ be one edge $uv$, let $T$ have one edge $t_1 t_2$, and let $W(t_1) = W(t_2) = G$. Then $W(T_1) \cap W(T_2) = G$. The second claim holds as printed, and $W(e)$ below uses only vertices, so it is not affected. Every graph of tree-width $\le w$ has a linked tree-decomposition of width $\le w$ in which no edge is in two of the graphs $W(t)$ (see the proof of (11.2)).} We define $W(e) = V(W(T_1) \cap W(T_2))$.

We say that a tree-decomposition $(T, W)$ of $G$ is \textit{linked} if for all $e_1, e_2 \in E(T)$ such that $|W(e_1)| = |W(e_2)| = k$ say, and $|W(e)| \ge k$ for all edges $e$ of the path of $T$ between $e_1$ and $e_2$, there are $k$ mutually vertex-disjoint paths of $G$ between $W(e_1)$ and $W(e_2)$. We need the following lemma.

\textbf{11.2} \textit{Let $G$ be a graph of tree-width $\le w$. Then there is a linked tree-decomposition $(T, W)$ of $G$, with width $\le w$, such that every vertex of $T$ has degree $\le 3$.}
